%% kapa: exam=1 round=23 year=2012 period=1 form=A subject=civil_law subject_ko=민법 seq=20 number=20
%% source: Q-Net 감정평가사 기출문제 (cst00309) artlSeq=5052578 file=raw/1008/5052578/2012년 제23회 감정평가사 제1차시험 기출문제(1교시 A형).pdf
%% conversion: kordoc 4.24.51 -> tools/gen1.py (auto); review=proofread-sonnet; flags=none
\begin{question}{20}
법률행위에 관한 설명으로 옳지 않은 것은? (다툼이 있으면 판례에 의함)
\begin{choices}
  \item 착오에 의한 의사표시의 취소권은 형성권이다.
  \item 복대리에 대해서는 표현대리가 성립되지 않는다.
  \item 대리인을 통한 계약체결시 의사표시의 하자 여부는 대리인을 표준으로 결정한다.
  \item 대리인이 다수인 때에는 특별한 사정이 없는 한 각자가 본인을 대리한다.
  \item 취소권자가 상대방으로부터 담보를 제공받은 경우 법정추인이 된다.
\end{choices}
\end{question}
\answer{2}
